% ID: SEL-D0082-OB-E03
% ESTADO: cotejada. Segunda pasada visual y borrador compilado revisados.
% Fuente DOC: Selectividad/sel_2012_matematicas_aplicadas/4 MATEMÁT. APLIC. CC. SS/EXAMEN 4-ANDALUCÍA.doc
% SHA-256 original: 8446c5da6b4b7cfc0f48b89d1b2d052ae0c8a1c4b5bc0b86b7e63399ce8656f8
% Página convertida: 2; opción B; ejercicio 3.

En una localidad hay solamente dos supermercados A y B. El $58\%$ de los habitantes compra en el A, el $35\%$ en el B y el $12\%$ compra en ambos.

Si se elige un ciudadano al azar, calcule la probabilidad de que:
\begin{enumerate}[label=\alph*)]
\item \textbf{(0.75 puntos)} Compre en algún supermercado.
\item \textbf{(0.5 puntos)} No compre en ningún supermercado.
\item \textbf{(0.5 puntos)} Compre solamente en un supermercado.
\item \textbf{(0.75 puntos)} Compre en el supermercado A, sabiendo que no compra en B.
\end{enumerate}
